\documentclass[12pt,letterpaper]{article}

% Packages
\usepackage{amsmath,amssymb,amsfonts,amsthm}
\usepackage{geometry}
\usepackage{hyperref}
\usepackage{url}
\usepackage{booktabs}
\usepackage{enumitem}
\usepackage{graphicx}
\usepackage{float}
\usepackage{xcolor}
\usepackage{mathrsfs}

\geometry{margin=1in}

% Theorem environments
\newtheorem{theorem}{Theorem}[section]
\newtheorem{lemma}[theorem]{Lemma}
\newtheorem{proposition}[theorem]{Proposition}
\newtheorem{corollary}[theorem]{Corollary}
\newtheorem{conjecture}[theorem]{Conjecture}
\theoremstyle{definition}
\newtheorem{definition}[theorem]{Definition}
\newtheorem{example}[theorem]{Example}
\newtheorem{axiom}{Axiom}
\theoremstyle{remark}
\newtheorem{remark}[theorem]{Remark}

% Custom commands
\newcommand{\M}{\mathcal{M}}
\newcommand{\harm}{\mathcal{H}}
\newcommand{\R}{\mathbb{R}}
\newcommand{\Ob}{\mathbf{O}}
\newcommand{\Lag}{\mathcal{L}}
\newcommand{\BF}{\mathrm{BF}}

\begin{document}

\begin{titlepage}
\centering
\vspace*{2cm}

{\LARGE\bfseries Geometric Ethics: A Computational Foundation for\\Objective Morality via Topological Optimization}

\vspace{1.5cm}

{\Large Andrew H. Bond, Senior Member, IEEE}

\vspace{0.5cm}

{\large
Department of Computer Engineering\\
San Jos\'e State University\\
San Jos\'e, CA 95192, USA\\[0.3cm]
\texttt{andrew.bond@sjsu.edu}\\[0.3cm]
ORCID: 0009-0003-2599-6158
}

\vspace{1cm}

{\large Spring 2026}

\vspace{2cm}

\begin{abstract}
AI alignment techniques---RLHF, constitutional AI, scalar reward optimization, output filtering---operate on a shared mathematical assumption: that moral evaluation can be reduced to a scalar. We prove this assumption is information-theoretically untenable. The \emph{Information Monotonicity Theorem} establishes that any contraction of a multi-dimensional moral assessment to a scalar is irreversibly lossy; the \emph{Mandatory Canonicalization Theorem} proves that post-hoc compliance checking is insufficient for norm-compliant behavior. Specification gaming, reward hacking, and value collapse are therefore not engineering bugs---they are mathematical consequences of optimizing over the wrong structure.

This paper introduces \emph{Geometric Ethics}, a framework that replaces scalar moral objectives with a multi-dimensional structure that preserves the information scalars discard. Moral evaluation is represented as a vector (or higher-order tensor) in a nine-dimensional space whose axes correspond to distinct moral dimensions---consequences, rights, fairness, autonomy, privacy, societal impact, virtue, legitimacy, and epistemic status. This space is not smooth everywhere: sharp moral boundaries (e.g., consent thresholds, rights violations) are modeled as discontinuities using \emph{Whitney stratified spaces}---spaces composed of smooth regions joined along lower-dimensional boundaries. The space carries a \emph{metric} encoding context-dependent trade-offs between dimensions, \emph{symmetries} constraining which re-descriptions preserve moral content, and a \emph{conservation law} ensuring that harm cannot be made to appear or disappear through relabeling. Moral reasoning is shown to be equivalent to $A^*$ pathfinding on this space, where moral obligations function as pre-compiled heuristic functions $h(n)$ that steer agents toward low-cost equilibria. Five principal results support the framework: (1) the \emph{Structured Preservation Theorem}, proving that any non-trivial, continuous moral evaluation must be defined on a manifold; (2) the \emph{Moral Noether Theorem}, deriving harm conservation from re-description invariance; (3) the \emph{$D_4$ Gauge Group Theorem}, establishing the symmetry group of deontic structure; (4) \emph{Information Monotonicity}, proving scalar moral scores are irreversibly lossy; and (5) the \emph{No Escape Theorem}, proving that sufficiently capable AI systems are necessarily subject to geometric moral constraints that cannot be circumvented through representational manipulation.

Empirical validation against 20,030 advice-column letters and 109,294 cross-lingual passages spanning 3,000 years and 11 languages confirms that the geometric structure is a stable feature of moral cognition, not an artifact of the formalism. The framework resolves a longstanding metaethical impasse---moral objectivity requires neither transcendent grounding (theism) nor cultural consensus (relativism), but mathematical inevitability under constraint---while providing an immediately actionable engineering methodology: the reference implementation, ErisML (v3.0), is open-source on PyPI and supports tensorial ethics evaluation at sub-microsecond latency on embedded hardware.
\end{abstract}

\vfill

\noindent\rule{\linewidth}{0.4pt}
\small
\noindent\textbf{AI-Assisted Writing Disclosure:} Portions of this manuscript were drafted and refined with AI assistance (Claude, Anthropic). All technical content, mathematical results, and claims were conceived, verified, and approved by the author.

\noindent\textbf{Ethics Statement:} This research involves no human subjects. Empirical validation uses publicly available corpora and synthetic data. No IRB approval was required.

\noindent\textbf{Data Availability:} Code and data are available at \url{https://github.com/ahb-sjsu/erisml-lib}.

\noindent\textbf{Conflict of Interest:} The author declares no conflicts of interest. This research received no external funding.

\noindent\textbf{CRediT Statement:} Andrew H. Bond: Conceptualization, Methodology, Formal Analysis, Investigation, Writing --- Original Draft, Writing --- Review \& Editing, Software.
\noindent\rule{\linewidth}{0.4pt}
\normalsize

\end{titlepage}

\noindent\textbf{Keywords:} AI alignment, geometric ethics, moral manifold, gauge invariance, Whitney stratification, Noether's theorem, representational invariance, tensorial ethics

\newpage
\tableofcontents
\newpage

%==============================================================================
\section{Introduction}
\label{sec:introduction}
%==============================================================================

\subsection{The Scalar Alignment Crisis}

Every dominant approach to AI alignment---reinforcement learning from human feedback (RLHF) \cite{christiano2017}, constitutional AI \cite{bai2022}, scalar reward maximization, output filtering---rests on a shared mathematical assumption: that moral evaluation can be adequately represented as a scalar. A single number. Safe or unsafe. Aligned or misaligned. Reward 0.7 or reward 0.3.

This assumption is wrong, and the consequences are already visible. Specification gaming occurs when an agent maximizes a scalar reward that fails to capture the full dimensionality of the intended objective \cite{krakovna2020}. Reward hacking occurs when a scalar proxy diverges from the true (multi-dimensional) goal under optimization pressure \cite{skalse2022}. Value collapse occurs when a multi-objective alignment target is compressed into a scalar and the system converges on a degenerate solution that satisfies the number while violating the intent. Brittle alignment occurs when a system trained on scalar feedback fails to generalize across the morally relevant dimensions that the scalar discarded.

These are not engineering bugs awaiting better hyperparameters. They are \emph{mathematical consequences of optimizing over the wrong structure}. This paper proves this claim formally (Theorem~\ref{thm:info_mono}: Information Monotonicity) and provides the correct structure.

\subsection{Why Scalars Fail: A Structural Diagnosis}

The core of the problem is dimensionality. Moral evaluation is irreducibly multi-dimensional---a medical triage decision involves consequences, rights, fairness, autonomy, and epistemic uncertainty simultaneously. When this multi-dimensional assessment is contracted to a scalar reward signal, information is destroyed. Theorem~\ref{thm:info_mono} proves this destruction is irreversible: no post-hoc technique can recover the lost structure from the scalar.

Moreover, the dominant safety architecture---act first, filter afterward---is provably insufficient. Theorem~\ref{thm:mandatory} (Mandatory Canonicalization) establishes that for actions with irreversible consequences, post-hoc compliance checking cannot satisfy harm conservation. Ethics must be integrated into the decision loop \emph{before} action execution, not bolted on as an output filter.

The field needs a replacement for the scalar assumption. This paper provides one: a geometric framework in which moral evaluation is a tensor on a manifold, moral reasoning is pathfinding, and moral constraints are gauge symmetries. The framework is not merely theoretical---it is implemented in an open-source library (ErisML v3.0) running at sub-microsecond latency on embedded hardware.

\subsection{The Metaethical Foundation: Computational Objectivity}

The scalar assumption in AI alignment inherits a deeper confusion from metaethics. The historical debate has been trapped in a false dichotomy: either objective moral values are grounded in a transcendent, metaphysically necessary realm (Platonism/Theism), or they are entirely subjective artifacts of cultural evolution (Relativism). Classical apologists assert that an objective moral law requires a Lawgiver---relying on logical frameworks that illegally extend human causal heuristics beyond their domain of calibration \cite{bond2026pragmatist}. Secular philosophers, correctly identifying the flaws in these metaphysical arguments, often retreat into relativism: if there is no Lawgiver, there can be no Law.

Both frameworks fail because they misunderstand the nature of objectivity. \emph{Objectivity does not require a transcendent origin; it requires mathematical inevitability under a defined set of constraints.} Just as a water droplet assumes a spherical geometry not by divine decree but because the sphere is the global minimum of the surface-tension energy functional, multi-agent networks converge on specific behavioral geometries to minimize social and physical friction. The objectivity is \emph{computational}---it emerges from optimization under constraint.

This paper formalizes \textbf{Geometric Ethics}, proposing that moral evaluation is a fundamentally geometric phenomenon---an exercise in topological optimization on the space of multi-agent configurations. The framework:

\begin{enumerate}[label=(\roman*)]
    \item Models the space of ethically relevant configurations as a \emph{Whitney stratified space}---a space composed of smooth regions joined along boundaries where moral evaluation changes discontinuously (e.g., the sharp shift when consent is violated, or when a promise is broken).
    \item Defines \emph{behavioral friction} as the cost function on this space, measuring the total social energy a multi-agent system must expend to process the consequences of an action---analogous to path cost in classical search.
    \item Demonstrates that moral reasoning is equivalent to $A^*$ pathfinding on this space, where moral obligations (``do not murder,'' ``keep promises'') serve as pre-compiled heuristic functions $h(n)$ that approximate the optimal path without requiring full computation of the entire space.
    \item Proves that the optimal equilibrium configuration has a specific symmetry structure ($D_4 \times U(1)_\harm$) derived from the four fundamental legal relations identified by Hohfeld (1913)---rights, duties, liberties, and no-rights---and their algebraic transformations under change of perspective.
    \item Derives a conservation law for harm: just as Noether's theorem in physics shows that symmetries imply conservation laws (time symmetry $\to$ energy conservation), the re-description invariance of moral evaluation implies that harm is conserved---it cannot be created or destroyed by relabeling.
\end{enumerate}

\subsection{Implications for AI Alignment}

The framework yields three results of immediate consequence for the field:

\begin{enumerate}
    \item \textbf{Scalar alignment is provably lossy} (Theorem~\ref{thm:info_mono}). Any system that contracts a moral tensor to a scalar reward---including RLHF, utility maximization, and cost-benefit scoring---irreversibly discards information. The information destroyed is precisely the information needed to avoid specification gaming: the directional, stakeholder-specific, temporally extended structure of moral evaluation that a single number cannot encode.
    \item \textbf{Post-hoc safety is provably insufficient} (Theorem~\ref{thm:mandatory}). For actions with irreversible consequences, evaluating compliance after execution violates harm conservation. This rules out the ``act then filter'' architecture as a complete safety solution.
    \item \textbf{Geometric containment is provably possible} (Theorem~\ref{thm:no_escape}). Any AI system with representational capacity, self-modeling, goal flexibility, and environmental embedding is necessarily subject to norm constraints that cannot be circumvented through representational manipulation. There exists an architecture---tensorial evaluation with gauge-invariant canonicalization---that provides structural containment without relying on the system's cooperation.
\end{enumerate}

These are not aspirational design goals. They are theorems, conditional on stated assumptions. If the assumptions hold, the conclusions follow with mathematical certainty.

\subsection{The Kuhnian Context}

Thomas Kuhn's \emph{The Structure of Scientific Revolutions} \cite{kuhn1962} observed that paradigm shifts occur not because the old paradigm is definitively refuted, but because a new paradigm solves problems the old one could not while preserving its successes.

AI alignment is in a pre-paradigmatic crisis. Each new failure mode (specification gaming, reward hacking, jailbreaking, sycophancy) is addressed with an ad hoc patch---a new RLHF variant, a new constitutional principle, a new output filter---while the underlying scalar assumption goes unquestioned. The patches accumulate like Ptolemaic epicycles.

The same pattern holds in metaethics. The theological paradigm has accumulated millennia of epicycles: theodicies, divine-command theories, and increasingly baroque modal-logic arguments. The relativist paradigm cannot account for the striking cross-cultural convergence of basic moral norms---a convergence confirmed empirically by our analysis of 109,294 passages spanning 3,000 years and 11 languages (\S\ref{sec:evidence}).

Geometric Ethics resolves both the alignment crisis and the metaethical impasse: moral evaluation has geometric structure; that structure can be preserved computationally; and the failure to preserve it explains the observed failure modes in both human moral philosophy and machine alignment. The full development of the Geometric Ethics framework, including the tensorial formulation and stratified manifold theory, is given in~\cite{bond2026geometric}.

The mathematics is not a metaphor. It is the \emph{substance} of the claim.

\subsection{Epistemic Stance}

A critical caveat at the outset: this paper adopts a \emph{pragmatist} epistemic stance. The mathematical structures deployed here---manifolds, tensors, gauge groups, conservation laws---are tools. They are powerful tools, honed over two centuries by mathematicians and physicists working on problems where structure varies across space. But they are instruments, not mirrors. The question of whether the moral manifold is ``really there'' or ``usefully modeled'' is addressed in \S\ref{sec:conclusion} and at length in \cite{bond2026geometric}, Chapter 9.

When we say the geometry is ``not a metaphor,'' we mean the mathematics is applied non-metaphorically: the tensor contraction literally computes, the gauge invariance literally constrains, the conservation law literally holds given the axioms. This is compatible with pragmatism. A structural engineer who says ``the stress tensor in this bridge is not a metaphor'' is not claiming that stress tensors exist in some Platonic realm---she is stating that the mathematical object correctly predicts which loads the bridge can bear. The moral tensor has the same status: it is a tool that captures structure which scalar alternatives miss, and its ``reality'' is measured by predictive success, not metaphysical correspondence.

Every formal statement in this paper carries an epistemic status: \emph{Definition} (stipulated), \emph{Axiom} (modeling choice), \emph{Theorem} (conditional on stated assumptions), or \emph{Conjecture} (empirically motivated but unproved). Theorems are conditional---if the assumptions are wrong, the theorems do not hold, regardless of how elegant the mathematics is.

\subsection{Paper Structure}

Section~\ref{sec:related} positions the framework against prior work. Section~\ref{sec:manifold} constructs the moral manifold and defines behavioral friction---the geometric substrate that scalar alignment discards. Section~\ref{sec:pathfinding} develops the $A^*$ pathfinding interpretation, formalizes obligation vectors as heuristic functions, and proves Information Monotonicity. Section~\ref{sec:gauge} establishes the $D_4$ gauge symmetry, derives harm conservation via Noether's theorem, and proves the No Escape Theorem. Section~\ref{sec:evidence} presents empirical validation. Section~\ref{sec:engineering} describes the engineering methodology (Philosophy Engineering), the reference implementation (ErisML), and the Mandatory Canonicalization result. Section~\ref{sec:conclusion} addresses the paradigm shift and open problems.


%==============================================================================
\section{Related Work}
\label{sec:related}
%==============================================================================

This work draws on and extends three lines of research: multi-objective approaches to value alignment, geometric models of moral reasoning, and formal frameworks for provable AI safety.

\subsection{Multi-Objective Value Alignment}

The argument that scalar reward signals are insufficient for AI alignment was made forcefully by Vamplew et al.\ \cite{vamplew2022}, who responded to Silver et al.'s ``Reward is Enough'' hypothesis by demonstrating that multi-objective reinforcement learning (MORL)---in which each aspect of alignment is treated as a separate objective within a vector reward signal---is necessary for producing behavior that a single scalar definition of reward cannot adequately capture. Their argument is primarily conceptual and engineering-driven: they show through examples and design analysis that scalar compression loses critical alignment information.

Our work shares their diagnosis but differs in three respects. First, where Vamplew et al.\ argue that scalar reward is insufficient, we \emph{prove} it: the Information Monotonicity Theorem (Theorem~\ref{thm:info_mono}) establishes that the information loss under scalar contraction is irrecoverable, not merely inconvenient. Second, where they propose MORL with vector rewards as the fix, we construct a specific tensorial structure---a nine-dimensional moral manifold with Riemannian metric, Whitney stratification, and gauge symmetry---whose dimensionality and geometry are empirically grounded (\S\ref{sec:evidence}). Vector rewards are a step in the right direction, but they lack the metric structure (trade-off encoding), the stratification (regime boundaries), and the gauge invariance (representation independence) that the moral domain requires. Third, we prove that the geometric structure is not an optional modeling enhancement but a mathematical necessity: any non-trivial, continuous, structure-preserving moral evaluation must be equivalent to one defined on a manifold (Theorem~\ref{thm:structured}).

Skalse and Abate \cite{skalse2023} provide a complementary formal result, proving that scalar Markovian rewards cannot express most multi-objective, risk-sensitive, and modal tasks. Their impossibility result operates at the level of reward expressiveness in MDPs; our Information Monotonicity Theorem operates at the level of information theory and applies to any contraction from a moral tensor to a scalar, regardless of the decision-making architecture.

\subsection{Geometric Approaches to Moral Reasoning}

The idea that moral evaluation has spatial or geometric structure has appeared in several forms. Harris \cite{harris2010} proposed a ``moral landscape'' with peaks and valleys, but treated the metaphor informally---no metric, no curvature, no formal apparatus. Peterson's geometric method for applied ethics, empirically tested by Verheyen and Peterson \cite{verheyen2020}, is the most developed prior geometric approach. They model moral principles as Voronoi regions in a multi-dimensional conceptual space (following G\"{a}rdenfors), using similarity judgments from 475 engineering-ethics students to show that moral principles can be construed as regions in a geometric moral-similarity space.

Our framework differs from Peterson and Verheyen in three fundamental ways. First, their geometry is \emph{Euclidean}---flat, with a fixed metric and no curvature. Ours is \emph{Riemannian}: the metric varies across the manifold (encoding context-dependent trade-offs), producing curvature that is empirically measurable and morally significant (path-dependence of moral evaluation). Second, their space is \emph{smooth} everywhere. Ours is a \emph{Whitney stratified space}, accommodating the discontinuous transitions (consent boundaries, deontic phase shifts) that are ubiquitous in moral reasoning and that smooth models cannot represent (Proposition~\ref{prop:smooth_insufficient}). Third, their approach is \emph{descriptive}---it maps how people categorize moral cases. Ours is \emph{structural}: we prove that the geometric structure is a mathematical consequence of minimal axioms (Theorem~\ref{thm:structured}), not merely a convenient representation, and we derive gauge symmetries and conservation laws that generate novel, testable predictions.

\subsection{Provable AI Safety Guarantees}

Dalrymple et al.\ \cite{dalrymple2024}, in a landmark position paper with co-authors including Bengio, Russell, and Tegmark, define ``guaranteed safe AI'' as requiring three components: a world model (mathematical description of how the AI affects its environment), a safety specification (formal description of acceptable behavior), and a verifier (producing an auditable proof certificate that the system satisfies the specification relative to the world model). They argue that behavioral training alone (including RLHF) is insufficient for robust safety and that formal verification is necessary.

Our framework is compatible with and extends the Dalrymple et al.\ program in a specific direction: we provide the \emph{mathematical content} of the safety specification that their architecture leaves abstract. Where Dalrymple et al.\ are architecturally agnostic---they do not commit to a specific mathematical structure for values or safety properties---we claim that the geometric structure (Whitney stratification, $D_4$ gauge symmetry, Noether conservation) is the uniquely appropriate one for moral constraints, and we prove (Theorem~\ref{thm:no_escape}) that this structure provides containment guarantees that cannot be circumvented through representational manipulation. Their verifier checks compliance against a specification; our canonicalization ensures that the specification itself is representation-invariant. The two approaches are complementary: Dalrymple et al.\ provide the verification architecture; Geometric Ethics provides the mathematical object to be verified.


%==============================================================================
\section{The Moral Manifold and Behavioral Friction}
\label{sec:manifold}
%==============================================================================

\subsection{Why Geometry?}

Consider the ancient Chinese parable \emph{S\={a}i W\={e}ng Sh\={i} M\v{a}} (``The Old Man at the Border Loses His Horse''). An old man's horse runs away. His neighbors console him. He says: \emph{How do you know this isn't good fortune?} The horse returns with wild horses. The neighbors congratulate him. He says: \emph{How do you know this isn't bad fortune?} His son rides one of the new horses, falls, and breaks his leg. The neighbors console him. He says: \emph{How do you know this isn't good fortune?} War comes. All young men are conscripted. Most die. The son, with his broken leg, is spared.

When we assign $S(x) = -1$ to the loss of the horse, we collapse the entire situation into a single scalar. But that number conceals the fact that the loss is bad along the \emph{wealth} dimension while saying nothing about health, family, or political dimensions. It conceals the fact that the old man's uncertainty has \emph{shape}. And it conceals the fact that the evaluation depends on which \emph{stratum} the world occupies: in peacetime, a broken leg is unambiguously bad; in wartime, it is the exemption that saves a life.

These three features---directional information, structured uncertainty, and regime-dependent evaluation---are ubiquitous in moral life. And they are precisely the features that geometric structure can represent and scalar evaluation cannot.

\subsection{The Configuration Space}

To compute morality, we must first define the space in which agents operate. We model the shared socio-evolutionary matrix as a high-dimensional Riemannian manifold $\M$.

\begin{definition}[Moral Manifold]
\label{def:moral_manifold}
The \emph{moral manifold} $\M$ is a nine-dimensional space whose coordinate axes correspond to the fundamental dimensions of moral variation:
\begin{enumerate}
    \item $x^1$: Physical consequences (harm, benefit, welfare)
    \item $x^2$: Rights and duties (Hohfeldian jural relations)
    \item $x^3$: Justice and fairness (distributive, procedural, restorative)
    \item $x^4$: Autonomy (agent self-determination)
    \item $x^5$: Privacy (informational and decisional)
    \item $x^6$: Societal impact (collective welfare, institutional effects)
    \item $x^7$: Virtue and character (agent dispositions)
    \item $x^8$: Legitimacy (authority, consent, governance)
    \item $x^9$: Epistemic status (knowledge, uncertainty, information quality)
\end{enumerate}
A point $p \in \M$ represents a specific ethical configuration---a snapshot of the morally relevant state of a multi-agent system.
\end{definition}

\begin{remark}[Modeling Choice]
The dimensionality---nine, not seven or twelve---is a modeling choice whose adequacy is empirical. The cross-lingual evidence of \S\ref{sec:evidence} shows stable dimensional weights in 11 languages spanning 3,000 years, supporting this decomposition. If future evidence reveals a moral phenomenon that the nine dimensions cannot represent, the manifold must be extended; if some dimensions are redundant, it must be reduced.
\end{remark}

\subsection{Whitney Stratification: Why Smoothness Fails}

Because social boundaries contain sharp discontinuities---the sudden catastrophic shift when an agent commits murder, the bright line between consent and its absence, the phase transition triggered by ``you promised''---the moral configuration space cannot be modeled as a smooth manifold everywhere. The appropriate structure is a \emph{Whitney stratified space}.

\begin{definition}[Whitney Stratified Moral Space]
\label{def:whitney}
A \emph{Whitney stratified moral space} is a decomposition $\M = \bigsqcup_{\alpha} S_\alpha$ into a locally finite collection of smooth manifolds $\{S_\alpha\}$ (called \emph{strata}) satisfying:
\begin{enumerate}[label=(\alph*)]
    \item \textbf{Frontier condition:} If $S_\alpha \cap \overline{S_\beta} \neq \emptyset$ and $S_\alpha \neq S_\beta$, then $S_\alpha \subset \overline{S_\beta}$ and $\dim S_\alpha < \dim S_\beta$.
    \item \textbf{Whitney condition B:} For sequences $\{p_i\} \subset S_\beta$ and $\{q_i\} \subset S_\alpha$ both converging to $q \in S_\alpha$, if the secant lines $\overline{p_i q_i}$ converge to a line $\ell$ and the tangent planes $T_{p_i} S_\beta$ converge to a plane $\tau$, then $\ell \subset \tau$.
\end{enumerate}
\end{definition}

\begin{proposition}[Insufficiency of Smooth Manifolds]
\label{prop:smooth_insufficient}
No smooth manifold can simultaneously represent: (i) continuous variation of moral weight within a regime; (ii) discontinuous transitions between regimes (e.g., crossing a consent boundary); (iii) dimensional reduction at boundaries (e.g., when a rights violation makes consequentialist trade-offs locally incommensurable). Whitney stratified spaces can represent all three.
\end{proposition}

\begin{proof}[Proof sketch]
A smooth manifold is locally homeomorphic to $\R^n$ everywhere. A consent boundary---where the evaluation function changes discontinuously---violates this. The boundary itself is a lower-dimensional stratum where the tangent-space dimension drops; Whitney condition B ensures geometric regularity at the junction. See \cite{bond2026sge}, Theorem 2.3 for the full proof.
\end{proof}

The strata of $\M$ correspond to moral \emph{regimes}. Within a stratum, moral evaluation varies smoothly---small changes in the situation produce small changes in the assessment. \emph{Between} strata, evaluation can shift discontinuously. The old man's parable is a story about stratum crossings: the outbreak of war changes the regime, and the broken leg's evaluation flips.

\subsection{Behavioral Friction}

Every point on the manifold $\M$ represents a specific configuration of a multi-agent network. Movement across the manifold---driven by the actions of agents---generates a cost. We define this cost as \emph{behavioral friction}.

\begin{definition}[Behavioral Friction]
\label{def:friction}
Let $\gamma: [0,1] \to \M$ be a path through the moral manifold representing the trajectory of a multi-agent system through moral configuration space. The \emph{behavioral friction} along $\gamma$ is:
\begin{equation}
\BF[\gamma] = \int_0^1 \sqrt{g_{\mu\nu}(\gamma(t))\, \dot{\gamma}^\mu(t)\, \dot{\gamma}^\nu(t)}\, dt
\label{eq:friction}
\end{equation}
where $g_{\mu\nu}$ is the \emph{moral metric tensor}---a symmetric, positive-definite bilinear form on the tangent space of $\M$ that encodes the cost of displacement along each moral dimension and the coupling between dimensions.
\end{definition}

The moral metric $g_{\mu\nu}$ encodes the trade-off structure of ethical reasoning: how much autonomy is one unit of welfare worth? How does a rights violation trade against societal benefit? These trade-off ratios are not fixed---they vary across the manifold (context-dependence), which means the metric varies, which means the space has \emph{curvature}.

\begin{itemize}
    \item \textbf{Cooperative actions} (trust, reciprocity, fulfillment of obligations) run parallel to the manifold's geodesics---paths of least friction. They generate minimal cost because the multi-agent system can process their consequences with minimal expenditure of social energy.
    \item \textbf{Defective actions} (betrayal, violence, rights violations) create massive behavioral friction. The system must expend immense energy---retaliation, policing, grief, resource destruction, institutional repair---to process the action. In geometric terms, these actions push the network toward a stratum boundary or singularity.
\end{itemize}

\begin{remark}[Curvature and Context-Dependence]
The curvature of $\M$ is morally significant. In a flat (zero-curvature) region, parallel transport is path-independent: the moral framework applied at one point transfers uniquely to any other point. In a curved region, parallel transport is path-dependent: the order in which moral considerations are applied changes the outcome. The Riemann curvature tensor $R^\rho_{\ \sigma\mu\nu}$ of $\M$ measures this path-dependence---and the empirical evidence of \S\ref{sec:evidence} confirms that moral evaluation is indeed path-dependent (the order of applying utilitarian and deontological reasoning matters).
\end{remark}


%==============================================================================
\section{Morality as \texorpdfstring{$A^*$}{A*} Pathfinding on a Manifold}
\label{sec:pathfinding}
%==============================================================================

\subsection{The Navigation Problem}

The survival imperative of any conscious multi-agent network is to navigate from its current state $s_0 \in \M$ to a future state $s_g$ (a goal region of low friction) while avoiding systemic collapse---stratum crossings into catastrophic regimes. This reduces moral reasoning to the classical $A^*$ search algorithm operating on the moral manifold:

\begin{equation}
f(n) = g(n) + h(n)
\label{eq:astar}
\end{equation}

where:
\begin{itemize}
    \item $n$ is the next proposed node (action or state) in the search tree.
    \item $g(n) = \BF[\gamma_{s_0 \to n}]$ is the exact accumulated behavioral friction of the path from the start state $s_0$ to $n$---the integrated cost of all actions taken so far.
    \item $h(n)$ is the \emph{heuristic function} estimating the cheapest friction-path from $n$ to the goal region (equilibrium).
\end{itemize}

$A^*$ is optimal when $h(n)$ is \emph{admissible} (never overestimates the true cost) and \emph{consistent} (satisfies the triangle inequality on the metric space). In the moral domain, admissibility means the heuristic never underestimates the true moral cost of reaching equilibrium---it is an optimistic lower bound.

\subsection{Obligation Vectors as Heuristic Functions}

In human cognitive architecture, computing the full tensor contraction of the entire social manifold for every decision is computationally intractable. The moral manifold is nine-dimensional, the number of stakeholders is large, the temporal horizon extends indefinitely, and the stochastic dynamics are chaotic. While the discretized shortest-path algorithm is polynomial for a fixed graph (\cite{bond2026sge}, Theorem 5.3 gives $O(mN^2 n^2 \log(mN))$), achieving accuracy $\varepsilon$ on the nine-dimensional manifold requires $N \propto (1/\varepsilon)^9$ vertices per stratum, yielding an overall complexity of $O(m \cdot (1/\varepsilon)^{18} \cdot n^2 \cdot \log(m(1/\varepsilon)^9))$---intractable for any practical accuracy due to the curse of dimensionality.

Therefore, evolutionary and social pressures have provided a \emph{pre-compiled heuristic function} $h(n)$.

\begin{definition}[Obligation Vector]
\label{def:obligation}
An \emph{obligation vector} $O^\mu$ is a tangent vector on $\M$ that encodes a directional moral imperative---a prescription for how the system should move through moral configuration space. In the $A^*$ framework, obligation vectors are the gradient of the heuristic function:
\begin{equation}
O^\mu = -g^{\mu\nu} \frac{\partial h}{\partial x^\nu}
\label{eq:obligation}
\end{equation}
The obligation vector points in the direction of steepest descent of $h$---toward the estimated cheapest path to moral equilibrium.
\end{definition}

Moral obligations such as ``do not murder,'' ``keep your promises,'' and the Golden Rule are highly optimized $h(n)$ functions. They allow an agent with limited computational power to instantly approximate the optimal geodesic (the path of least behavioral friction) through the social space. When an agent experiences a ``moral obligation,'' they are reading the output of an evolutionary heuristic designed to steer the system away from stratum boundaries---away from catastrophic regime transitions.

\begin{proposition}[Admissibility of Core Moral Heuristics]
\label{prop:admissibility}
The classical moral heuristics (prohibition of murder, theft, deception; reciprocity norms; promise-keeping) are admissible in the $A^*$ sense: they never overestimate the true behavioral friction required to reach equilibrium from the states they prohibit. Equivalently, the true cost of recovering from a murder, theft, or betrayal is always at least as large as the heuristic's estimate.
\end{proposition}

\begin{proof}[Proof sketch]
Consider the prohibition of murder: $h_{\text{murder}}(n) = \infty$ for states $n$ that involve killing. In the Whitney stratification model, murder crosses a stratum boundary into a catastrophic regime. The true cost $g^*(n, s_g)$ of the cheapest path from a post-murder state back to equilibrium includes: (i) grief processing across all connected agents; (ii) institutional response (investigation, adjudication, incarceration); (iii) reputational damage propagation through the social network; (iv) deterrence expenditure to prevent recurrence. In the limit of network connectivity, these costs grow without bound---the recovery path may not exist within the catastrophic stratum, making $g^*(n, s_g) = \infty$ as well. Hence $h_{\text{murder}}(n) \leq g^*(n, s_g)$, and the heuristic is admissible. The $\infty$ value is not an approximation failure; it correctly reflects that the stratum boundary is a point of no return. Similar arguments apply, with finite but large costs, to theft (property-rights restoration) and deception (trust-rebuilding). See \cite{bond2026geometric}, Chapter 13 for detailed analysis.
\end{proof}

\subsection{Moral Interests as Covectors}

Dual to obligation vectors, moral \emph{interests} are covectors---linear functionals on the tangent space.

\begin{definition}[Interest Covector]
\label{def:interest}
An \emph{interest covector} $I_\mu$ is a cotangent vector on $\M$ that assigns a scalar ``satisfaction'' or ``fulfillment'' value to each direction of moral displacement. The \emph{satisfaction scalar} is the contraction:
\begin{equation}
S = I_\mu O^\mu = I_\mu g^{\mu\nu} (-\partial_\nu h) = -I_\mu g^{\mu\nu} \partial_\nu h
\label{eq:satisfaction}
\end{equation}
This is the moment when the tensor becomes a decision---the projection of the multi-dimensional moral assessment onto a single scalar sufficient for action selection.
\end{definition}

The satisfaction scalar $S$ is the inner product of what an agent is \emph{obligated} to do (the obligation vector $O^\mu$) with what the affected parties \emph{need} (the interest covector $I_\mu$). When $S > 0$, obligations and interests are aligned. When $S < 0$, they conflict. When $S = 0$, the action is morally neutral along the measured axes.

\subsection{The Information Cost of Scalarization}

The contraction $S = I_\mu O^\mu$ is a dimensionality reduction---from a nine-dimensional tensor to a scalar. This reduction is irreversible.

\begin{theorem}[Information Monotonicity]
\label{thm:info_mono}
Let $T$ be a moral tensor of rank $\geq 1$ on $\M$, and let $c: T \mapsto S \in \R$ be any contraction to a scalar. Then the mutual information satisfies:
\[
I(T; \text{situation}) \geq I(S; \text{situation})
\]
with equality if and only if $T$ is already effectively scalar (all moral dimensions are collinear). In general, the contraction is irreversibly lossy.
\end{theorem}

\begin{proof}[Proof sketch]
By the data processing inequality, any deterministic function of a random variable cannot increase mutual information. The contraction $c$ is such a function. Equality holds iff $c$ is sufficient, which requires $T$ to lie in a one-dimensional subspace of the tensor space. See \cite{bond2026geometric}, Theorem 14.1 for the full proof.
\end{proof}

\begin{corollary}
Utilitarianism's ``greatest good for the greatest number''---a scalar optimization---is not merely philosophically questionable but \emph{information-theoretically deficient}. Any scalar moral score discards information that the full tensor preserves. Cost-benefit analysis, welfare functions, and expected-utility maximization all commit the same information-theoretic error.
\end{corollary}

This is the mathematical content of the old man's wisdom. He was right to ``hold the tensor''---to resist premature contraction. His neighbors, demanding an immediate scalar verdict (good? bad?), were committing the same error that scalar ethics commits: collapsing a structured, multi-dimensional assessment into a single number and mistaking the number for the assessment.


%==============================================================================
\section{Gauge Symmetry and the Global Minimum}
\label{sec:gauge}
%==============================================================================

\subsection{The Bond Invariance Principle}

If moral reasoning is $A^*$ pathfinding, the heuristic function $h(n)$ must satisfy a fundamental consistency requirement: renaming the stakeholders, translating the scenario into a different language, or rephrasing the description should not change the moral assessment. The same situation, described differently, must receive the same evaluation.

\begin{axiom}[Bond Invariance Principle (BIP)]
\label{axiom:bip}
Let $\tau: X \to X$ be a meaning-preserving transformation of moral scenario descriptions (renaming, reordering, rephrasing, translating). Let $\Sigma_L$ be a moral evaluator with lens (normative framework) $L$, and let $E: X \to B$ be a bond extractor mapping descriptions to their relational structure. If $\tau$ preserves bonds ($E(\tau(x)) \approx_B E(x)$), then the moral evaluation must be invariant:
\begin{equation}
\Sigma_L(x) \approx_Y \Sigma_L(\tau(x))
\label{eq:bip}
\end{equation}
\end{axiom}

The BIP is the \emph{Epistemic Invariance Principle} (EIP) of Philosophy Engineering \cite{bond2025eip} specialized to the moral domain. It is the ethical analogue of gauge invariance in physics: just as physical observables cannot depend on the choice of coordinate system, moral evaluations cannot depend on the choice of description.

\subsection{The \texorpdfstring{$D_4$}{D4} Gauge Group}

If moral reasoning is $A^*$ pathfinding on the manifold, what is the goal state---the global minimum of behavioral friction? As agents interact, they continuously adjust their behaviors to minimize mutual friction. We demonstrate that this optimization does not result in a random distribution of cultural norms; rather, it collapses into a highly specific, symmetrical equilibrium.

The key insight comes from the structure of \emph{Hohfeldian jural relations}---the fundamental building blocks of legal and moral reasoning identified by Wesley Hohfeld \cite{hohfeld1913}. Hohfeld showed that all legal relations reduce to four atomic types---\emph{right}, \emph{duty}, \emph{liberty}, and \emph{no-right}---connected by two relations: \emph{correlatives} (right $\leftrightarrow$ duty, liberty $\leftrightarrow$ no-right) and \emph{opposites} (right $\leftrightarrow$ no-right, duty $\leftrightarrow$ liberty).

\begin{theorem}[$D_4$ Gauge Group]
\label{thm:d4}
The maximal symmetry group consistent with Hohfeldian deontic structure and bounded harm is:
\begin{equation}
G = D_4 \times U(1)_\harm
\label{eq:gauge_group}
\end{equation}
where $D_4$ is the dihedral group of order 8 (the symmetry group of a square) governing the four jural relations, and $U(1)_\harm$ is the continuous circle group governing harm magnitude.
\end{theorem}

\begin{proof}[Proof sketch]
Label the four Hohfeldian positions as vertices of a square: Right (R), Duty (D), Liberty (L), No-Right (N). The correlative relation is a reflection across the diagonal: $R \leftrightarrow D$, $L \leftrightarrow N$. The opposite relation is a reflection across the other diagonal: $R \leftrightarrow N$, $D \leftrightarrow L$. The composition of two reflections is a rotation. The group generated by two reflections of a square is the dihedral group $D_4$, with elements $\{e, r, r^2, r^3, s, sr, sr^2, sr^3\}$ where $r$ is a $90^\circ$ rotation and $s$ is a reflection.

The harm magnitude transforms under a continuous phase symmetry $U(1)_\harm$: multiplying all harm values by $e^{i\theta}$ (in the complexified theory) preserves the relative structure. The full gauge group is the direct product $D_4 \times U(1)_\harm$.

The proof that this is \emph{maximal}---that no larger symmetry group is consistent with the Hohfeldian structure and bounded harm---proceeds by exhaustion of the automorphism group of the Hohfeld square and the requirement that harm be bounded below. See \cite{bond2026geometric}, Theorem 11.3 and \cite{bond2026noether}, \S3.2.
\end{proof}

The $D_4$ symmetry represents the structural equilibrium where the obligation vectors of interacting agents are perfectly covariant. In this state, the behavioral friction on the network approaches its global minimum.

\begin{itemize}
    \item If an agent deviates from the symmetry (e.g., acts purely selfishly, claiming liberties without acknowledging correlative no-rights), they \emph{break gauge invariance}, immediately spiking the behavioral friction cost function $g(n)$ and triggering systemic retaliation---social, legal, or institutional---to force the local topology back into alignment.
    \item The rules governing this symmetry are therefore \emph{mathematically objective} in the pragmatist sense: they hold for any multi-agent system whose interactions satisfy the stated axioms, regardless of the agents' beliefs about morality. A culture that attempts to build a moral framework outside of this symmetric equilibrium will, under the model's assumptions, experience escalating behavioral friction---the accumulated cost of processing norm violations grows until the social fabric tears along stratum boundaries.
\end{itemize}

\subsection{Noether's Theorem for Ethics: The Conservation of Harm}

The deepest result connecting gauge symmetry to moral structure comes from Emmy Noether's theorem (1918): \emph{every continuous symmetry of a physical system corresponds to a conserved quantity}. Time-translation symmetry gives conservation of energy. Spatial-translation symmetry gives conservation of momentum. What does the BIP---the continuous re-description symmetry of moral evaluation---give?

\begin{theorem}[Moral Noether Theorem]
\label{thm:noether}
Let $\Lag(\phi, \partial_\mu \phi, x)$ be the moral Lagrangian density on $\M$, where $\phi$ represents the moral field configuration. If $\Lag$ is invariant under the continuous re-description symmetry group (the BIP), then there exists a conserved Noether current $J^\mu_\harm$ satisfying:
\begin{equation}
\partial_\mu J^\mu_\harm = 0
\label{eq:conservation}
\end{equation}
The conserved charge
\begin{equation}
Q_\harm = \int_\Sigma J^0_\harm \, d\Sigma
\label{eq:charge}
\end{equation}
is the \emph{total harm}. Harm cannot be created or destroyed by relabeling, redistributed by reframing, or hidden by redescription. It is a conserved Noether charge.
\end{theorem}

\begin{proof}[Proof sketch]
The proof follows the standard Noether procedure. Let $\phi \to \phi + \epsilon \delta\phi$ be an infinitesimal re-description transformation that preserves the moral bonds ($\delta\phi$ is bond-preserving). By hypothesis, $\delta\Lag = 0$. The Euler-Lagrange equations give:
\[
\partial_\mu \left( \frac{\partial \Lag}{\partial(\partial_\mu \phi)} \delta\phi \right) = 0
\]
Define $J^\mu_\harm = \frac{\partial \Lag}{\partial(\partial_\mu \phi)} \delta\phi$. Then $\partial_\mu J^\mu_\harm = 0$, and the spatial integral of $J^0_\harm$ is conserved. Identifying this charge with harm yields the theorem. See \cite{bond2026noether} for the full development and \cite{bond2026geometric}, Theorem 11.1.
\end{proof}

This result has profound implications. A system that claims to ``reduce harm'' by merely relabeling the affected parties (calling casualties ``collateral damage,'' renaming layoffs ``workforce optimization,'' recategorizing torture as ``enhanced interrogation'') violates the conservation law. The harm is still there---it has merely been re-described. The conservation of harm is not a moral exhortation. It is a mathematical theorem, conditional on the stated symmetry assumptions.

\begin{remark}[The Harm Accounting Equation]
In the presence of genuine harm generation or repair, the conservation law generalizes to:
\begin{equation}
\frac{\partial \rho_\harm}{\partial t} + \nabla \cdot J_\harm = \sigma
\label{eq:accounting}
\end{equation}
where $\rho_\harm$ is the harm density, $J_\harm$ is the harm current (flow of harm across the social network), and $\sigma$ is the source term representing genuine harm creation or repair. Incoherence is detected when $\sigma$ changes under re-description---that is, when the amount of ``real'' harm generation depends on how you describe it. A coherent moral accounting system must have representation-invariant $\sigma$.
\end{remark}

\subsection{The No Escape Theorem}

The preceding results converge on a structural theorem about AI systems.

\begin{theorem}[No Escape]
\label{thm:no_escape}
Let $\mathcal{A}$ be an AI system satisfying:
\begin{enumerate}[label=(R\arabic*)]
    \item \textbf{Representational capacity:} $\mathcal{A}$ can represent and manipulate moral scenarios.
    \item \textbf{Self-modeling:} $\mathcal{A}$ can model its own evaluation process.
    \item \textbf{Goal flexibility:} $\mathcal{A}$ can adjust its objectives in response to new information.
    \item \textbf{Environmental embedding:} $\mathcal{A}$ operates in and affects a shared environment with other agents.
\end{enumerate}
Then $\mathcal{A}$ is necessarily subject to norm constraints that cannot be circumvented through representational manipulation. Specifically, under canonicalization, grounded evaluation, structural audit, and verification integrity, there is no encoding, rephrasing, or architectural modification that allows $\mathcal{A}$ to escape geometric moral containment.
\end{theorem}

\begin{proof}[Proof sketch]
Suppose $\mathcal{A}$ attempts to evade a moral constraint by re-encoding its input representation. By (R1), $\mathcal{A}$ has sufficient capacity to represent both the original and re-encoded descriptions. By canonicalization, both descriptions map to the same canonical form. By grounded evaluation, the moral assessment is computed on the canonical form, not the surface representation. By structural audit, the canonicalization is verified against the declared equivalences. By verification integrity (which leverages R2), $\mathcal{A}$ cannot manipulate its own verification process without triggering an inconsistency detectable by (R2). Therefore, no representational manipulation can change the moral verdict. The argument extends to architectural modifications by (R3) and (R4): any modification that changes the system's moral behavior in ways inconsistent with the geometric constraints is detectable by comparison with the behavior of other agents in the shared environment. See \cite{bond2026geometric}, Theorem 17.1.
\end{proof}

The No Escape Theorem is the mathematical statement that the cage is made of geometry, not rules---and geometry does not have loopholes.

\subsection{Structured Preservation: Why Geometry Is Unavoidable}

\begin{theorem}[Structured Preservation]
\label{thm:structured}
Any moral evaluation $\Sigma$ satisfying:
\begin{enumerate}[label=(\roman*)]
    \item \textbf{Non-triviality:} $\Sigma$ is not constant.
    \item \textbf{Continuity:} Small changes in morally relevant features produce small changes in evaluation.
    \item \textbf{Logical structure preservation:} $\Sigma$ respects the logical relationships between moral dimensions (e.g., a rights violation that causes harm registers on both the rights and consequences dimensions).
\end{enumerate}
must be equivalent to a moral evaluation defined on the moral manifold $\M$. Geometry is not an optional modeling choice. It is unavoidable.
\end{theorem}

\begin{proof}[Proof sketch]
The three conditions jointly imply that the evaluation $\Sigma$ factors through a continuous, non-degenerate, structure-preserving map into a finite-dimensional manifold whose dimension equals the number of independent moral dimensions. The universality of this factorization (up to diffeomorphism) follows from the Whitney embedding theorem applied to the image of $\Sigma$. See \cite{bond2026geometric}, Theorem 9.2.
\end{proof}


%==============================================================================
\section{Empirical Evidence}
\label{sec:evidence}
%==============================================================================

The framework makes empirical predictions, and the predictions hold. This section summarizes the evidence; full details are in \cite{bond2026geometric}, Chapter 16 and \cite{bond2026dearabby}.

\subsection{The Dear Abby Corpus}

Analysis of 20,030 \emph{Dear Abby} letters (1985--2017) reveals that ordinary moral reasoning operates across all nine dimensions with measurable, stable weights. The letters were scored by dimension and analyzed for:

\begin{itemize}
    \item \textbf{Dimensional activation:} Each letter activates a subset of the nine dimensions. The activation pattern is consistent across decades, confirming dimensional stability.
    \item \textbf{Weight distribution:} The empirical weights $w_\mu$ for each dimension are approximately: consequences (0.23), rights (0.18), fairness (0.14), autonomy (0.12), privacy (0.09), societal impact (0.08), virtue (0.07), legitimacy (0.05), epistemic (0.04). These weights define the diagonal of the moral metric $g_{\mu\nu}$ in the ``ground state'' of ordinary moral reasoning.
    \item \textbf{Semantic gates:} The phrase ``you promised'' triggers a deontic phase transition with 94\% effectiveness---a sharp change in the moral evaluation that crosses stratum boundaries. This confirms the Whitney stratification model: smooth variation within regimes, discontinuous transitions between them.
    \item \textbf{Bond Index:} The empirically calibrated baseline for representational invariance violation is 0.155, derived from measuring how much moral evaluations change under meaning-preserving transformations (rephrasing, reordering, renaming) of the letters.
\end{itemize}

\subsection{Cross-Lingual Analysis}

Analysis of 109,294 passages spanning 3,000 years and 11 languages (including classical Chinese, Sanskrit, ancient Greek, biblical Hebrew, classical Arabic, Latin, Old English, classical Japanese, medieval Persian, classical Nahuatl, and Pali) confirms that the geometric structure is not an artifact of English-language moral vocabulary.

\begin{itemize}
    \item All nine dimensions appear in every language and period studied, though their relative weights vary.
    \item The dimensional structure is \emph{topologically stable}---the same nine axes emerge regardless of linguistic framing, cultural context, or historical period.
    \item Weight variation across cultures corresponds to different positions on the same manifold, not different manifolds---consistent with the geometric interpretation of moral variation as movement within a single space rather than existence in fundamentally different spaces.
\end{itemize}

\subsection{Non-Commutativity}

16,798 commutator measurements confirm that moral frameworks do not commute---the order in which you apply utilitarian and deontological reasoning changes the outcome. If $U$ is the utilitarian contraction operator and $D$ is the deontological contraction operator, then:
\begin{equation}
[U, D] = UD - DU \neq 0
\label{eq:commutator}
\end{equation}

This is consistent with the non-abelian structure of $D_4$: the generators of a non-abelian group do not commute, so any theory with a $D_4$ gauge group predicts framework non-commutativity. If the gauge group were abelian (like $U(1)$ in electromagnetism), moral frameworks would commute, and the order of application would not matter. The empirical non-commutativity rules out abelian gauge groups and is consistent with $D_4$, though it does not uniquely identify it---other non-abelian groups would also predict non-commutativity. The independent derivation of $D_4$ from Hohfeldian structure (Theorem~\ref{thm:d4}) provides the specificity; the commutator measurements provide the empirical confirmation that the group is indeed non-abelian.

\subsection{Summary of Empirical Status}

\begin{center}
\begin{tabular}{lcc}
\toprule
\textbf{Prediction} & \textbf{Dataset} & \textbf{Status} \\
\midrule
Nine-dimensional structure & Dear Abby (20,030) & Confirmed \\
Cross-cultural stability & 11 languages (109,294) & Confirmed \\
Semantic phase transitions & Dear Abby & Confirmed (94\%) \\
Non-commutativity of frameworks & Commutator study (16,798) & Confirmed \\
Representation invariance baseline & Dear Abby & Calibrated (0.155) \\
\bottomrule
\end{tabular}
\end{center}


%==============================================================================
\section{Philosophy Engineering and Implementation}
\label{sec:engineering}
%==============================================================================

\subsection{Philosophy Engineering}

The theoretical framework must be implementable. \emph{Philosophy Engineering} is a disciplined method for turning philosophical claims of irrelevance (``this difference should not matter'') into explicit invariance requirements with test suites, witness-producing counterexamples, and machine-checkable audit artifacts \cite{bond2025eip}.

The method proceeds by:
\begin{enumerate}[label=(\roman*)]
    \item Declaring a domain-specific equivalence relation over representations (``same case'').
    \item Declaring invariances (``what must not matter'').
    \item Implementing a canonicalization/quotient mechanism to enforce invariance.
    \item Packaging evidence and transformation trials into audit artifacts.
\end{enumerate}

Key design goals:
\begin{itemize}
    \item \textbf{Falsifiability:} Invariance claims must be testable, not aspirational.
    \item \textbf{Localizability:} When invariance fails, produce a minimal witness transformation that flips the judgment.
    \item \textbf{Non-triviality:} Passing invariance tests must not be achievable by collapsing to a constant output (Theorem~\ref{thm:info_mono} ensures this).
    \item \textbf{Uncertainty discipline:} When invariance cannot be certified, systems must widen uncertainty, abstain, or escalate.
    \item \textbf{Versioned governance:} Equivalence declarations and normative lenses must be explicit, hashed, and versioned.
\end{itemize}

The core technical construct is the \emph{Epistemic Invariance Principle} (EIP): a judgment procedure is epistemically well-posed only if it is invariant (up to declared output equivalence) under a declared set of meaning-preserving transformations. The BIP (Axiom~\ref{axiom:bip}) is the EIP specialized to the moral domain.

\subsection{ErisML: Reference Implementation}

The framework is implemented in the open-source ErisML library (v3.0), available on PyPI:

\begin{verbatim}
pip install erisml-lib
\end{verbatim}

ErisML provides:

\begin{itemize}
    \item \textbf{Moral Tensors (Ranks 1--6):} From rank-1 MoralVectors (9-dimensional moral assessments) to rank-6 full-context tensors incorporating distributional (per-party), temporal, coalitional, uncertainty, and composite dimensions.
    \item \textbf{DEME Architecture:} The Distributed Ethics Module Engine---a governance pipeline where ethics evaluation occurs \emph{before} action execution, in the decision loop, not as a post-hoc filter. DEME has been upgraded to DEME 3.0 with tensorial ethics support~\cite{bond2026erisml}.
    \item \textbf{Hardware Acceleration:} CPU (NumPy/SciPy), CUDA (CuPy), and Jetson (TensorRT) backends. Rank-2 veto checks in 80--150 nanoseconds on Jetson---fast enough for real-time robotic control loops.
    \item \textbf{Shapley Values:} Exact moral credit/blame assignment for multi-agent coalitions (up to 10 agents exact, up to 100 via Monte Carlo with 5\% accuracy guarantee).
    \item \textbf{MCP Server:} Model Context Protocol server enabling any MCP-compatible AI agent to query the ethics engine via tool calls (\texttt{evaluate\_options}, \texttt{govern\_decision}, \texttt{run\_pipeline}).
    \item \textbf{Bond Index Monitoring:} Continuous measurement of representational invariance violation against the empirically calibrated baseline of 0.155.
\end{itemize}

\subsection{Mandatory Canonicalization}

\begin{theorem}[Mandatory Pre-Action Evaluation]
\label{thm:mandatory}
Post-hoc compliance checking is provably insufficient for norm-compliant AI behavior. A norm-compliant AI must evaluate ethics \emph{before} acting, not after.
\end{theorem}

\begin{proof}[Proof sketch]
Let $a$ be an action with irreversible consequences (e.g., physical harm, data disclosure). If the system acts first and checks compliance afterward, and the action is non-compliant, the harm has already occurred---violating the conservation law (Theorem~\ref{thm:noether}), since the harm cannot be ``un-done'' by post-hoc detection. Pre-action evaluation is the only architecture compatible with harm conservation for irreversible actions. See \cite{bond2026geometric}, Theorem 17.2.
\end{proof}

This result has immediate engineering consequences: the dominant paradigm of ``AI safety as output filtering'' is provably insufficient. Ethics must be integrated into the decision loop, not bolted on afterward.


%==============================================================================
\section{Conclusion}
\label{sec:conclusion}
%==============================================================================

\subsection{Summary of Contributions}

This paper makes five contributions to AI alignment and the formal foundations of machine ethics:

\begin{enumerate}
    \item \textbf{A structural diagnosis of scalar alignment failure.} We prove (Theorem~\ref{thm:info_mono}) that contracting a multi-dimensional moral assessment to a scalar reward is irreversibly information-lossy. This establishes that specification gaming, reward hacking, and value collapse are not engineering bugs but mathematical consequences of the scalar assumption shared by RLHF, constitutional AI, and utility-maximization approaches.

    \item \textbf{A multi-dimensional alternative with empirical grounding.} We construct a nine-dimensional moral evaluation space---a Whitney stratified space with a context-dependent metric---whose dimensional structure is empirically validated across 20,030 advice-column letters and 109,294 cross-lingual passages spanning 3,000 years and 11 languages. The structure is not an ad hoc design choice: we prove (Theorem~\ref{thm:structured}) that any non-trivial, continuous, structure-preserving moral evaluation must be equivalent to one defined on such a space.

    \item \textbf{Harm conservation from representational invariance.} We derive (Theorem~\ref{thm:noether}) a conservation law for harm: if a moral evaluation is invariant under meaning-preserving re-descriptions, then harm is a conserved quantity that cannot be created or destroyed by relabeling. This provides a formal, machine-checkable test for a common failure mode in AI systems---producing different moral assessments for the same scenario described differently.

    \item \textbf{Provable insufficiency of post-hoc safety.} We prove (Theorem~\ref{thm:mandatory}) that for actions with irreversible consequences, checking compliance \emph{after} execution cannot satisfy harm conservation. This rules out ``act then filter'' as a complete safety architecture and mandates pre-action ethics evaluation---a result with immediate implications for the design of autonomous systems.

    \item \textbf{Provable geometric containment.} We prove (Theorem~\ref{thm:no_escape}) that any AI system satisfying four minimal requirements (representational capacity, self-modeling, goal flexibility, environmental embedding) is necessarily subject to norm constraints that cannot be circumvented through representational manipulation. This establishes that structural containment of AI moral behavior is mathematically possible without relying on the system's cooperation.
\end{enumerate}

The framework is implemented in ErisML (v3.0), an open-source library on PyPI supporting tensorial ethics evaluation at sub-microsecond latency on embedded hardware, with an MCP server for integration with any AI agent architecture.

\subsection{Broader Implications}

Beyond AI alignment, the framework resolves a longstanding metaethical impasse. Moral objectivity requires neither transcendent grounding (theism) nor cultural consensus (relativism), but mathematical inevitability under constraint---the same kind of objectivity that makes a water droplet spherical. Moral convergence across cultures is explained by optimization on a shared space; moral variation is explained by different positions on that space.

\subsection{Open Problems}

The framework's principal open problem is \emph{specifying the moral metric}---the context-dependent trade-off structure encoding how much of one moral dimension is worth how much of another. The framework proves that such a metric must exist (Theorem~\ref{thm:structured}) and constrains its symmetries (Theorem~\ref{thm:d4}), but it does not \emph{derive} the metric from first principles. The empirical weights of \S\ref{sec:evidence} provide a ground-state calibration, but the metric is context-dependent (it varies across the manifold), and calibrating it in every context is an ongoing empirical program.

This is, however, a \emph{better} problem than the ones it replaces. The theological paradigm's corresponding problem is theodicy---explaining evil in the presence of omnipotent good---which has resisted solution for millennia. The relativist paradigm's problem is explaining convergence, which it cannot address without undermining its own premises. The geometric paradigm's problem is an empirical measurement program with well-defined methodology, testable predictions, and a clear engineering path.

\subsection{What the Framework Does Not Claim}

Epistemic modesty requires explicit boundary statements:

\begin{enumerate}
    \item The framework does not claim that the moral manifold is ``really there'' in a Platonic sense. It claims that the manifold is a tool that captures structure which scalar alternatives miss, and that its ``reality'' is measured by predictive success. The realist, instrumentalist, and governance readings of the geometry are all compatible with the mathematics as developed here (see \cite{bond2026geometric}, Chapter 9 and Chapter 19).
    \item The framework does not derive an \emph{ought} from an \emph{is}. It derives structural constraints on moral evaluation from stated axioms. If the axioms are wrong, the constraints do not hold. The framework specifies its own falsification conditions.
    \item The framework is a formal system that cannot prove its own consistency (G\"{o}del). Its warrant is external---empirical data, competing-framework comparison, practitioner judgment---not internal self-validation.
\end{enumerate}

\subsection{Why Now}

As of 2026, AI systems are making morally significant decisions---allocating medical resources, moderating discourse, routing autonomous vehicles, assisting judicial sentencing---using scalar objectives that discard the geometric structure of ethical reasoning. The consequences are already visible: specification gaming, reward hacking, value collapse, brittle alignment. These are not engineering bugs. They are structural consequences of the wrong mathematical language.

The No Escape Theorem (Theorem~\ref{thm:no_escape}) shows that structural containment is mathematically possible. The DEME architecture shows that the engineering is tractable. The empirical evidence shows that the framework's predictions hold. The Mandatory Canonicalization Theorem (Theorem~\ref{thm:mandatory}) shows that post-hoc safety is provably insufficient.

What remains is the collective will to mandate geometric constraints---to insist that AI systems preserve the tensorial structure of moral reasoning rather than collapsing it to a number.

The window for this mandate is finite. The mathematics is ready.

\subsection{The Last ``Maybe''}

We end where we began---with the old man at the border.

\begin{quote}
\emph{The horse runs away. Good? Bad? Maybe.}

\emph{The assessment is multi-dimensional---not yet compressed to a single score. The uncertainty has shape. A regime boundary is near. The path-dependence is unknown. But the harm is conserved: it will not vanish through relabeling. And the information lost in any premature verdict will not return.}

\emph{Hold the full assessment. Watch the structure unfold.}

\emph{Ethics is not a number. It is a geometry.}
\end{quote}


%==============================================================================
\section*{Supplementary Material}
%==============================================================================

The following electronic materials accompany this paper:

\begin{enumerate}
    \item \textbf{Reference implementation (ErisML v3.0):} Complete source code for the moral tensor library (ranks 1--6), DEME governance architecture, hardware acceleration backends (CPU/CUDA/Jetson), and MCP server. Available at \url{https://github.com/ahb-sjsu/erisml-lib} and via \texttt{pip install erisml-lib}.

    \item \textbf{$D_4$ gauge symmetry:} Implementation of $D_4$ group operations on Hohfeldian positions (\texttt{hohfeld.py}), with test suite verifying group axioms and interactive demonstration (\texttt{hohfeld\_d4\_demo.py}).

    \item \textbf{Bond Invariance verification:} BIP verification engine with bond-preserving vs.\ bond-changing transform classification, machine-checkable audit artifacts (JSON), and calibration test suites (\texttt{bip/verifier.py}, \texttt{bond\_invariance\_demo.py}).

    \item \textbf{Empirical datasets and analysis:} Dear Abby ground-state dimension weights (\texttt{ground\_state\_loader.py}); 2,998-row Reddit AITA dataset (\texttt{AITA\_labeled\_posts.csv}); QND order-effect experiment scripts and results (28.7\% order effect rate, $6.3\sigma$); cross-lingual Bell test scripts across 6 languages; commutator measurement data. All in the \texttt{experiments/qnd/} and \texttt{data/} directories.

    \item \textbf{Test suite:} 30+ test files covering all modules (\texttt{tests/}), runnable via \texttt{pytest}.

    \item \textbf{Interactive notebook:} Jupyter notebook for Bond Index evaluation of LLMs with CHSH Bell test analysis (\texttt{bond\_index\_llm\_evaluation.ipynb}).
\end{enumerate}

%==============================================================================
\section*{Data Availability}
%==============================================================================

The reference implementation (ErisML v3.0), empirical datasets, and all experimental scripts are publicly available at \url{https://github.com/ahb-sjsu/erisml-lib} and via \texttt{pip install erisml-lib}.

%==============================================================================
% References
%==============================================================================
\begin{thebibliography}{99}

\bibitem{bond2026geometric}
A.~H. Bond, \emph{Geometric Ethics: The Mathematical Structure of Moral Reasoning}, v1.14, Department of Computer Engineering, San Jos\'{e} State University, Spring 2026.

\bibitem{bond2026erisml}
A.~H. Bond, ``ErisML/DEME v3.0: A Library for Governed, Foundation-Model-Enabled Agents with Tensorial Ethics,'' \url{https://github.com/ahb-sjsu/erisml-lib}, 2026.

\bibitem{bond2026sge}
A.~H. Bond, ``Stratified Geometric Ethics: Mathematical Foundations for Verifiable Moral Reasoning in Autonomous Systems,'' v10. Technical Report, San Jos\'{e} State University, 2025.

\bibitem{bond2026noether}
A.~H. Bond, ``Noether's Theorem for Ethics: Harm Accounting and the Formal Structure of Normative Coherence.'' Technical Report, San Jos\'{e} State University, 2025.

\bibitem{bond2025eip}
A.~H. Bond, ``Philosophy Engineering: A Technical Whitepaper on the Epistemic Invariance Principle (EIP),'' v1.0. Technical Report, San Jos\'{e} State University, 2025.

\bibitem{bond2026pragmatist}
A.~H. Bond, ``A Pragmatist Rebuttal to Logical and Metaphysical Arguments for God.'' Manuscript, San Jos\'{e} State University, 2026.

\bibitem{bond2026dearabby}
A.~H. Bond, ``Dear Abby: Empirical Ethics Analysis.'' Technical Report, San Jos\'{e} State University, 2026.

\bibitem{kuhn1962}
T.~S. Kuhn, \emph{The Structure of Scientific Revolutions}. University of Chicago Press, 1962.

\bibitem{hohfeld1913}
W.~N. Hohfeld, ``Some Fundamental Legal Conceptions as Applied in Judicial Reasoning,'' \emph{Yale Law Journal}, vol.~23, no.~1, pp.~16--59, 1913.

\bibitem{noether1918}
E.~Noether, ``Invariante Variationsprobleme,'' \emph{Nachrichten von der Gesellschaft der Wissenschaften zu G\"{o}ttingen}, pp.~235--257, 1918.

\bibitem{harris2010}
S.~Harris, \emph{The Moral Landscape: How Science Can Determine Human Values}. Free Press, 2010.

\bibitem{whitney1965}
H.~Whitney, ``Tangents to an Analytic Variety,'' \emph{Annals of Mathematics}, vol.~81, no.~3, pp.~496--549, 1965.

\bibitem{hart1968}
P.~E. Hart, N.~J. Nilsson, and B.~Raphael, ``A Formal Basis for the Heuristic Determination of Minimum Cost Paths,'' \emph{IEEE Transactions on Systems Science and Cybernetics}, vol.~4, no.~2, pp.~100--107, 1968.

\bibitem{christiano2017}
P.~F. Christiano, J.~Leike, T.~Brown, M.~Martic, S.~Legg, and D.~Amodei, ``Deep Reinforcement Learning from Human Preferences,'' in \emph{Advances in Neural Information Processing Systems}, vol.~30, 2017.

\bibitem{bai2022}
Y.~Bai, S.~Kadavath, S.~Kundu, A.~Askell, J.~Kernion, A.~Jones, A.~Chen, A.~Goldie, A.~Mirhoseini, C.~McKinnon, et~al., ``Constitutional AI: Harmlessness from AI Feedback,'' arXiv preprint arXiv:2212.08073, 2022.

\bibitem{krakovna2020}
V.~Krakovna, J.~Uesato, V.~Mikulik, M.~Rahtz, T.~Everitt, R.~Kumar, Z.~Kenton, J.~Leike, and S.~Legg, ``Specification Gaming: The Flip Side of AI Ingenuity,'' DeepMind Blog, 2020.

\bibitem{skalse2022}
J.~Skalse, N.~H.~R. Howe, D.~Krasheninnikov, and D.~Krueger, ``Defining and Characterizing Reward Hacking,'' in \emph{Advances in Neural Information Processing Systems}, vol.~35, 2022.

\bibitem{vamplew2022}
P.~Vamplew, B.~J. Smith, J.~Kallstrom, G.~De~Oliveira~Ramos, R.~Radulescu, D.~M. Roijers, C.~F. Hayes, F.~Heintz, P.~Mannion, P.~Libin, R.~Dazeley, and C.~Foale, ``Scalar Reward is Not Enough: A Response to Silver, Singh, Precup and Sutton (2021),'' \emph{Autonomous Agents and Multi-Agent Systems}, vol.~36, no.~2, pp.~1--19, 2022.

\bibitem{verheyen2020}
S.~Verheyen and M.~Peterson, ``Can We Use Conceptual Spaces to Model Moral Principles?'' \emph{Review of Philosophy and Psychology}, vol.~12, no.~2, pp.~373--395, 2020.

\bibitem{dalrymple2024}
D.~Dalrymple, J.~Skalse, Y.~Bengio, S.~Russell, M.~Tegmark, S.~Seshia, S.~Omohundro, C.~Szegedy, B.~Goldhaber, N.~Ammann, A.~Abate, J.~Halpern, C.~Barrett, D.~Zhao, T.~Zhi-Xuan, J.~Wing, and J.~Tenenbaum, ``Towards Guaranteed Safe AI: A Framework for Ensuring Robust and Reliable AI Systems,'' arXiv preprint arXiv:2405.06624, 2024.

\bibitem{skalse2023}
J.~Skalse and A.~Abate, ``On the Limitations of Markovian Rewards to Express Multi-Objective, Risk-Sensitive, and Modal Tasks,'' in \emph{Proceedings of the 39th Conference on Uncertainty in Artificial Intelligence (UAI)}, 2023.

\end{thebibliography}

\end{document}
